\begin{frame}{자주 묻는 질문 (9.2)}
  \small
  \begin{itemize}
    \item \kb{Q. ReLU 미분이 0/1인데 왜 ``완전히'' 안 되나?}
      곱의 인자가 $\sigma'$ 하나뿐이 아님: (1) $W_l$ 크기가 나쁘면
      ReLU여도 소실/폭발 — 초기화의 역할, (2) $z<0$이면 미분 0이라
      죽은 뉴런은 영구히 기여 0 — ReLU 변형이 보완, (3) 학습 도중
      스케일이 다시 흐름 — BatchNorm의 역할.
      \textbf{``활성함수+초기화+정규화'' 세 라인이 함께} 작동해야
      해결.
    \item \kb{Q. 프레임워크 기본값이면 초기화 신경 안 써도 되나?}
      전형적 구조(완전연결 ReLU)에는 맞다 — PyTorch 기본값이
      Xavier/He. 하지만 (a) 스킵 커넥션/레지듀얼, (b) 커스텀 층
      (합성곱, LSTM 셀)에서는 ``fan-in''이 명확히 정의되지 않아
      기본값이 적용 안 됨. ``학습이 안 된다/발산한다''를 디버깅할
      때 \textbf{첫 질문은 항상 ``각 층의 그래디언트·활성 노름을
      재보라''}.
    \item \kb{Q. 출력층 활성함수는?} \textbf{문제 형태}가 결정:
      회귀 $\to$ 항등(MSE 짝), 이진분류 $\to$ 시그모이드(교차엔트로피
      짝), 다중분류 $\to$ softmax(Week 17 언어모델에서 재회).
    \item \kb{실습 노트북}: 10층 시그모이드 553$\times$ 감소, 세 초기화의
      분산 갈림, dying ReLU 97\%, 활성함수 6종 비교, 7층 XOR
      tanh+Xavier(NaN) vs ReLU+He(0.25) vs Adam(수렴).
  \end{itemize}
\end{frame}
